\section{Data Collection}\label{c:data}

\subsection{Method}

Collection was shared across the three team members. The procedure recorded in
\texttt{data/\allowbreak COLLECTION\_\allowbreak TEMPLATE.md} and followed for every
statement was:

\begin{enumerate}[leftmargin=1.4em, itemsep=0.05em]
  \item Listen in an ordinary setting --- taxis, chop houses, roadside stalls, bendskin
        ranks, the neighbourhood, the university.
  \item Write down \emph{the exact words spoken}, preserving slang, code-mixing, accents,
        incomplete sentences and mistakes. No cleaning up, no translation.
  \item Have a second Francanglais-familiar member check the transcription against what
        was heard, recording any disagreement rather than resolving it silently.
  \item Explain the coursework and obtain permission before retaining identifiable
        conversation; use anonymous identifiers and avoid unsafe collection settings.
\end{enumerate}

Statements were gathered over six days across a range of everyday situations rather than
in one sitting, so the corpus is not dominated by a single speaker, setting or topic.

\subsection{Coverage and attribution}

\CorpusTotal\ statements were collected and retained, above the 10--15 the assignment asks
for. Every one is attributed to the member who heard and transcribed it (\req{FR-4}).

\begin{center}\begin{tabular}{@{}llrr@{}}
\toprule
\textbf{Collector} & \textbf{Id} & \textbf{Statements} & \textbf{Accepted} \\
\midrule
Eyong Seanna Tabe & \texttt{ICTU20241297} & 7 & 7 \\
Kamdeu Yamdjeuson Neil Marshall & \texttt{ICTU20241386} & 8 & 6 \\
Tuheu Tchoubi Pempeme Moussa Fahdil & \texttt{ICTU20241393} & 7 & 7 \\
\midrule
\textbf{Total} & & \textbf{22} & \textbf{20} \\
\bottomrule
\end{tabular}
\end{center}

\subsection{Provenance enforced in software}

Authenticity is not left to good intentions --- it is a validated invariant.
\texttt{validate\_import} in \texttt{core/corpus.py} \emph{rejects} any \texttt{field}
record that does not also carry \texttt{manual\_transcription\_attested = true}
(\req{FR-3}). A statement cannot enter the graded corpus without an explicit human
attestation that it was listened to and transcribed by hand. All \CorpusAttested\
statements carry that attestation, and \texttt{tests/test\_field\_corpus.py} fails the
build if any of them loses it, if a statement becomes unattributable, or if a raw
transcription is silently altered (\req{FR-7}).

\subsection{Topic coverage}

The assignment lists ten themes. The analyzer models \TopicCount\ topics, matched by an
evidence-cited keyword table. A topic is always reported with the exact canonical words
that triggered it, so a reviewer can see why a label was attached; topic matching never
influences whether a statement is grammatically accepted (\req{FR-32}).

\begin{center}
\small
\begin{tabular}{@{}ll@{\hskip 26pt}ll@{}}
\toprule
\textbf{Required theme} & \textbf{Topic rule} & \textbf{Required theme} & \textbf{Topic rule} \\
\midrule
Taxi and commuting     & \term{commuting}          & Fuel scarcity          & \term{fuel} \\
Poor internet          & \term{internet}           & Roadside business      & \term{roadside\_business} \\
Limited electricity    & \term{electricity}        & Bendskin communication & \term{bendskin} \\
Market bargaining      & \term{market\_bargaining} & Security checkpoints   & \term{security} \\
Rainy season struggles & \term{rain}               & Life at the University & \term{university} \\
\midrule
\multicolumn{4}{@{}l@{}}{\small Added from observed vocabulary:
\term{food\_and\_drink}, \term{family}, \term{greetings}, \term{nightlife}, \term{housing}} \\
\bottomrule
\end{tabular}
\end{center}

The five additional topics were not planned. They were added because vocabulary kept
appearing that none of the ten prescribed themes covered, and discarding those statements
would have biased the corpus towards the themes we expected to find. The complete set of
raw transcriptions appears in Section~\ref{c:results}, paired with the verdict the grammar
returns for each.
